% Generated draft: requires mathematical and source review.
\documentclass[11pt]{article}
\usepackage[margin=25mm]{geometry}
\usepackage{fontspec}
\setmainfont{Latin Modern Roman}
\usepackage{amsmath,amssymb,mathtools,unicode-math}
\setmathfont{Latin Modern Math}
\usepackage{enumitem,needspace,xurl,hyperref}
\setlength{\parindent}{0pt}
\setlength{\parskip}{5pt}
\setlength{\emergencystretch}{3em}
\newcommand{\N}{\mathbb{N}}
\newcommand{\Z}{\mathbb{Z}}
\newcommand{\Q}{\mathbb{Q}}
\newcommand{\R}{\mathbb{R}}
\newcommand{\C}{\mathbb{C}}
\newcommand{\F}{\mathbb{F}}
\newcommand{\A}{\mathbb{A}}
\newcommand{\PP}{\mathbb{P}}
\newcommand{\E}{\mathbb{E}}
\newcommand{\eps}{\varepsilon}
\newcommand{\dd}{\,\mathrm{d}}
\begin{document}
\section*{UnsolvedMath ID: 30000557}
\textbf{Source problem code:} catalogue entry
\subsection{Definitions and mathematical statement}
Let \((M,\omega)\) be a closed symplectic manifold, with the symplectic form included in the notation \(M\). Let \(\operatorname{Ham}(M,\omega)\) be its group of smooth Hamiltonian time-one maps. For \(\phi\in\operatorname{Ham}(M,\omega)\), the Hofer norm is
\[
\|\phi\|_H=\rho_H(\operatorname{id},\phi)
=\inf_{\phi_H^1=\phi}\int_0^1
\bigl(\max_M H_t-\min_M H_t\bigr)\,\dd t,
\]
where the infimum is over smooth time-dependent Hamiltonians generating \(\phi\). Additive time-dependent constants in \(H\) do not change the displayed oscillation.

With commutator convention \([f,g]=fgf^{-1}g^{-1}\), define
\[
A(M,\omega)=\sup_{f,g\in\operatorname{Ham}(M,\omega)}\|[f,g]\|_H\in[0,\infty].
\]
The source asks whether this quantity is finite or infinite for a given closed symplectic manifold, and how the answer depends on the manifold. Finiteness means
\[
\exists C<\infty\quad\forall f,g\in\operatorname{Ham}(M,\omega),\qquad
\|[f,g]\|_H\leq C;
\]
infinity means that for every \(R>0\) some pair has \(\|[f,g]\|_H>R\). The constant may depend on \((M,\omega)\); no bound uniform across all symplectic manifolds is asked.

The supremum is over single commutators, not arbitrary products of commutators or the entire group's Hofer diameter. Nor is this a question about commutator length. The source specifies no manifold class beyond closed symplectic and gives no preferred finite or infinite answer. The supplied literature triage checked 2026-08-21 reports no authoritative resolution of the extracted general question, without a more specific theorem or example.
\subsection{Short English statement}
For a fixed closed symplectic manifold, are the Hofer norms of all single Hamiltonian commutators uniformly bounded, or can they be arbitrarily large? This is a commutator-subset diameter question, distinct from commutator length or the whole group's diameter.
\subsection{Sources}
\begin{itemize}
\item Local UnsolvedMath registry, permanent ID 30000557, catalogue entry, \emph{Commutator Diameters in Hamiltonian Groups}.
\item \emph{Geometric Group Theory, Hyperbolic Dynamics and Symplectic Geometry} (2007), Oberwolfach Reports, pp. 1991--2058, DOI \url{https://doi.org/10.4171/owr/2006/33}.
\item Supplied literature triage checked 2026-08-21 cites only that report; no separate general bound, unboundedness theorem, or class-specific example is supplied.
\end{itemize}
\end{document}
